\documentclass[11pt]{article}
\usepackage[margin=1in]{geometry}
\usepackage[T1]{fontenc}
\usepackage{lmodern,amsmath,amssymb,amsthm,microtype}
\usepackage[colorlinks=true,allcolors=blue]{hyperref}
\pdfinfoomitdate=1
\pdftrailerid{}
\ifdefined\pdfsuppressptexinfo\pdfsuppressptexinfo=15\fi
\newtheorem{theorem}{Theorem}
\newtheorem{lemma}[theorem]{Lemma}
\newtheorem{proposition}[theorem]{Proposition}
\newtheorem{corollary}[theorem]{Corollary}
\newtheorem{remark}[theorem]{Remark}
\newcommand{\F}{\mathbb F}
\newcommand{\R}{\mathbb R}
\newcommand{\diag}{\operatorname{diag}}
\newcommand{\gap}{\operatorname{gap}}
\newcommand{\wt}{\operatorname{wt}}
\newcommand{\Unif}{\operatorname{Unif}}
\providecommand{\SourceCommit}{SOURCE-NOT-BUILT}
\title{Product Marginals, Density Moments and Slow Spectra\\of Ising Transfer Kernels}
\author{Lluis Eriksson\\\small Independent researcher}
\date{13 November 2026}
\begin{document}
\maketitle
\begin{abstract}
Fix a positive Ising coupling. We construct strictly positive weights on
Boolean cubes whose labelled marginals on every set of at most a fixed
positive fraction of the coordinates equal the uniform product marginals,
while exponentially many normalized eigenvalues of the weighted transfer
kernel approach one exponentially fast in the number of coordinates.
The stationary law of the Perron Doob transform has exactly the same
marginals as well. The unweighted product kernel has a dimension-independent
gap. The construction combines binary codes with positive primal and dual
distances, a positive product admixture, and a finite block-norm estimate.
We prove the coding input by a direct entropy union bound and give explicit
finite spectral constants. As a negative control, we calculate the full
transfer spectrum of a single positive parity tilt: unlike its single-site
heat-bath gap, its transfer gap stays uniformly positive at fixed coupling.
Conversely, finite density moments of sufficiently high order give an
explicit dimension-independent transfer gap, without pointwise upper or
positive lower bounds, even allowing zero weights. Strictly positive
rare-spike densities with bounded lower-order moments can nevertheless
have exponentially collapsing gaps for this same operator. We give
explicit constants for both regimes; the interval between them and its
endpoints remain unresolved. The ingredients are classical and external
priority of these joint conclusions is not established. Global many-body
weights are allowed; no fixed-range Gibbs or local-dynamics claim follows.
\end{abstract}

\section{The question and the transfer convention}
Local marginal agreement and rapid relaxation are different kinds of
information. The distinction is familiar in limited independence and
Markov-chain theory, but a comparison must specify the dynamics: a slow
single-site heat-bath chain does not automatically give a slow transfer
operator. We make this distinction within a single fixed-coupling transfer
family and require exact, rather than approximate, agreement of marginal
input.

Let $\Omega_L=\F_2^L$, with Hamming distance $d(x,y)$. Fix $\beta>0$ and put
$q=e^{-2\beta}$. Our unweighted kernel is
\begin{equation}\label{eq:kernel}
 M_\beta(x,y)=q^{d(x,y)}
 =e^{-\beta L}\exp\!\left(\beta\sum_{i=1}^L(-1)^{x_i+y_i}\right).
\end{equation}
For a strictly positive probability law $\mu$ on $\Omega_L$, define
\begin{equation}\label{eq:weighted}
 A_\mu=\diag(\sqrt\mu)M_\beta\diag(\sqrt\mu).
\end{equation}
The kernel $M_\beta$ is a tensor power of
$\left(\begin{smallmatrix}1&q\\q&1\end{smallmatrix}\right)$ and is positive
definite. The same is true of $A_\mu$, which also has strictly positive
entries. Write its eigenvalues in decreasing order as
$\lambda_1>\lambda_2\ge\cdots>0$, counting multiplicities, and set
\[
 \gap(A_\mu)=1-\lambda_2/\lambda_1.
\]
Perron--Frobenius strictness gives a positive gap at each finite $L$; a
bound uniform in $L$ is a separate matter. Multiplying either kernel by
a positive scalar does not change any eigenvalue ratio.

Let $u_L=\Unif(\Omega_L)$. Tensorization gives
\begin{equation}\label{eq:product}
 \lambda_2(A_{u_L})/\lambda_1(A_{u_L})
       =\frac{1-q}{1+q}=\tanh\beta.
\end{equation}
For a law $\nu$, the literal marginal input of order $r$ is
\begin{equation}\label{eq:input}
 \mathcal M_r(\nu)=(\nu_S:S\subseteq[L],\ |S|\le r),
\end{equation}
where coordinate labels are retained. Full conditional distributions,
entrywise weight bounds, and the transfer matrix itself are not included
in this tuple.

\begin{theorem}\label{thm:main}
For each fixed $\beta>0$ there are constants $\alpha,R,\gamma>0$ such
that, for all sufficiently large integers $L$, a strictly positive law
$\mu_L$ on $\Omega_L$ has the following properties. Put
$r_L=\lfloor\alpha L\rfloor$ and $m_L=2^{\lceil RL\rceil}$.
\begin{enumerate}
\item $\mathcal M_{r_L}(\mu_L)=\mathcal M_{r_L}(u_L)$.
\item For every $2\le j\le m_L$,
\begin{equation}\label{eq:slow}
 0<1-\lambda_j(A_{\mu_L})/\lambda_1(A_{\mu_L})
                 \le 7e^{-\gamma L}.
\end{equation}
\item Let $v>0$ be the unit Euclidean Perron vector of $A_{\mu_L}$.
The reversible Perron Doob chain
\begin{equation}\label{eq:doob}
 P_L(x,y)=\frac{A_{\mu_L}(x,y)v_y}{\lambda_1 v_x},
 \qquad \pi_L(x)=v_x^2,
\end{equation}
has $\mathcal M_{r_L}(\pi_L)=\mathcal M_{r_L}(u_L)$ and the same
normalized spectrum as $A_{\mu_L}$.
\end{enumerate}
One admissible choice is
\begin{equation}\label{eq:parameters}
 R=\min\left\{\frac{1-h(1/4)}2,\frac{\beta}{4\log 2}\right\},
 \quad h(\alpha)<R/2,\quad
 \gamma=\beta/2-R\log 2\ge\beta/4,
\end{equation}
with $0<\alpha<1/2$ and binary entropy
$h(p)=-p\log_2p-(1-p)\log_2(1-p)$.
\end{theorem}

There are at least $m_L-1$ slow nonleading modes in this statement.
We do not assert that their eigenvalues coincide. The positivity and
strict gap at each finite size are compatible with their exponentially
small separation from the Perron root. The theorem does not impose a
uniform lower bound on the weights or a bounded interaction range.

\section{Codes with two distance conditions}
We include the coding argument to make all hypotheses of the spectral
construction explicit. Simultaneously good primal and dual codes are
classical; stronger algorithmic constructions are available
\cite{BrehmResch}. The relation between code projections and limited
independence is standard background for the marginal comparison
\cite{AGM}.

\begin{lemma}\label{lem:code}
Suppose $0<\alpha,\delta<1/2$ and
$h(\alpha)<R<1-h(\delta)$. Put $k=\lceil RL\rceil$.
If $k<L$ and
\begin{equation}\label{eq:codegate}
 2^{k-L(1-h(\delta))}+2^{-k+Lh(\alpha)}<1,
\end{equation}
there is a binary linear code $C\subseteq\F_2^L$ of dimension $k$
such that every nonzero codeword has weight at least $\delta L$ and
$C^\perp$ has no nonzero word of weight at most $\lfloor\alpha L\rfloor$.
Every coordinate projection of $C$ onto at most $\lfloor\alpha L\rfloor$
coordinates is surjective. The displayed conditions hold for all
sufficiently large $L$.
\end{lemma}
\begin{proof}
Choose a uniformly random binary $k\times L$ matrix $G$. For each fixed
nonzero row vector $z\in\F_2^k$, the word $zG$ is uniform in $\F_2^L$.
The probability that some such $zG$ has weight less than $\delta L$ is
at most
\[
 2^k2^{-L}\sum_{j\le\delta L}\binom Lj
       \le 2^{k-L(1-h(\delta))}.
\]
Failure of full row rank is included in this event, since it produces
a zero word from a nonzero message. For a fixed nonzero column vector
$w\in\F_2^L$, the vector $Gw$ is uniform in $\F_2^k$. Thus the probability
of a nonzero dual word with weight at most $\lfloor\alpha L\rfloor$ is
at most
\[
 2^{-k}\sum_{1\le j\le\alpha L}\binom Lj
       \le 2^{-k+Lh(\alpha)}.
\]
Their sum is less than one, so both failures can be avoided simultaneously.

For the entropy estimate, if $0<p<1/2$, set $z=p/(1-p)<1$. Since
$z^j\ge z^{pL}$ for $j\le pL$, the binomial theorem gives
\[
 \sum_{j\le pL}\binom Lj\le z^{-pL}(1+z)^L=2^{Lh(p)}.
\]
This bound applies to the inclusive sums and hence to the strict-weight
failure event as well. No integrality assumption on $\delta L$ is needed.

If the columns indexed by $S$ were dependent, their relation would give
a nonzero word of $C^\perp$ supported in $S$. Consequently $G_S$ has
rank $|S|$ whenever $|S|\le\lfloor\alpha L\rfloor$. Uniform messages
therefore give uniform projections onto $\F_2^S$. Finally, the exponents
in~\eqref{eq:codegate} tend to minus infinity, since $k=RL+O(1)$ and
both entropy inequalities are strict.
\end{proof}

\section{A finite positive concentration estimate}
The spectral step needs only the size and minimum distance of the code.
The dual distance will enter through the marginal statement.

\begin{proposition}\label{prop:finite}
Let $C\subseteq\F_2^L$ be a binary linear code with size $m\ge2$ and
minimum nonzero weight $d_C$. Suppose its coordinate projections onto
sets of size at most $r$ are surjective. For $0<\eta\le1/2$, let
\begin{equation}\label{eq:mixture}
 \mu=(1-\eta)\Unif(C)+\eta u_L,
 \qquad s=(m-1)e^{-2\beta d_C}.
\end{equation}
Then $\mu$ is strictly positive, $\mathcal M_r(\mu)=\mathcal M_r(u_L)$,
and every $2\le j\le m$ satisfies
\begin{equation}\label{eq:finite}
 0<1-\lambda_j(A_\mu)/\lambda_1(A_\mu)
       \le 2s+2m(\sqrt\eta+\eta).
\end{equation}
\end{proposition}
\begin{proof}
Uniform codewords have uniform projections by surjectivity, and a mixture
of two laws with the same marginals has those same marginals. The positive
product part gives strict positivity at every configuration.

Every codeword has probability
\[
 a=(1-\eta)/m+\eta/2^L\ge1/(2m).
\]
The principal block of $A_\mu$ indexed by $C$ is $aM_{CC}$. Its diagonal
is $a$, and the off-diagonal row sums of $M_{CC}$ are at most $s$:
distinct codewords differ by a nonzero codeword of weight at least $d_C$.
Thus all eigenvalues of $M_{CC}$ lie in $[1-s,1+s]$. For example, the
norm of its symmetric off-diagonal part is bounded by the maximum
absolute row sum. Principal interlacing gives
\begin{equation}\label{eq:interlace}
 \lambda_j(A_\mu)\ge\lambda_m(A_\mu)\ge a(1-s)
       \quad(1\le j\le m).
\end{equation}
Independently, a coordinate Rayleigh vector gives $\lambda_1(A_\mu)\ge a$.

Write $A_\mu$ in $C,C^c$ blocks. The outside principal block is positive
semidefinite, and its trace is
$\sum_{C^c}\mu=\eta(1-m/2^L)\le\eta$, so its norm is at most $\eta$.
For the rectangular cross block $B$, all entries of $M_\beta$ are at most
one, and therefore
\[
 \|B\|^2\le\|B\|_{\rm F}^2
 \le\left(\sum_C\mu\right)\left(\sum_{C^c}\mu\right)\le\eta.
\]
The symmetric block matrix with off-diagonal blocks $B,B^*$ and zero
diagonal has norm $\|B\|$. Consequently
\[
 \bigl\|A_\mu-\diag(aM_{CC},0)\bigr\|\le\sqrt\eta+\eta,
 \qquad
 \lambda_1(A_\mu)\le a(1+s)+\sqrt\eta+\eta.
\]
Subtract~\eqref{eq:interlace}, divide by $\lambda_1\ge a$, and use
$a\ge1/(2m)$ to obtain~\eqref{eq:finite}. Perron strictness proves the
strict lower inequality. The upper bound is valid even when it is
larger than one; no small-size success is claimed.
\end{proof}

\begin{proof}[Proof of Theorem~\ref{thm:main}, parts one and two]
Use~\eqref{eq:parameters} and Lemma~\ref{lem:code} with $\delta=1/4$.
The choice of $R$ satisfies $0<R<1-h(1/4)$ and $R\log2\le\beta/4$.
Choose codes for sufficiently large $L$, and put
\[
 m=2^{\lceil RL\rceil},\qquad \eta=e^{-2\beta L}/m^2.
\]
This number lies in $(0,1/4]$. Since $d_C\ge L/4$ and
$m\le2e^{R(\log2)L}$,
\[
 2(m-1)e^{-2\beta d_C}\le4e^{-(\beta/2-R\log2)L}
          =4e^{-\gamma L}.
\]
Also, using $m\ge2$ and $\gamma\le\beta$,
\[
 2m(\sqrt\eta+\eta)
   =2e^{-\beta L}+2e^{-2\beta L}/m
   \le3e^{-\beta L}\le3e^{-\gamma L}.
\]
Proposition~\ref{prop:finite} now proves both claims. The condition $k<L$
ensures $m<2^L$, so every asserted eigenvalue index exists.
\end{proof}

\section{Stationary marginals survive the Perron transform}
The weighting law in a transfer congruence need not equal the stationary
law of its associated Markov chain. Matching the former alone would leave
a significant interpretive gap. Code symmetry pays the latter exactly.

\begin{lemma}\label{lem:symmetry}
For the mixture~\eqref{eq:mixture}, the Perron stationary law
$\pi(x)=v_x^2$ satisfies $\mathcal M_r(\pi)=\mathcal M_r(u_L)$.
\end{lemma}
\begin{proof}
For every $c\in C$, translation $x\mapsto x+c$ preserves $\mu$, Hamming
distance, and hence $A_\mu$. The corresponding permutation commutes with
$A_\mu$ and takes a positive unit Perron vector to another such vector.
Uniqueness forces $v_{x+c}=v_x$, so $\pi$ is invariant under translations
by $C$. If $|S|\le r$, every translation on $\F_2^S$ lifts to some
$c\in C$ by the assumed surjectivity. The $S$-marginal of $\pi$ is
therefore invariant under all translations and must be uniform.
\end{proof}

\begin{proof}[Completion of Theorem~\ref{thm:main}]
The Perron eigenvector equation gives row sums one for~\eqref{eq:doob}.
Symmetry gives detailed balance:
\[
 \pi(x)P_L(x,y)=v_xA_\mu(x,y)v_y/\lambda_1
              =\pi(y)P_L(y,x).
\]
Conjugation by $\diag(v)$ shows that $P_L$ is similar to $A_\mu/\lambda_1$.
It has strictly positive entries and the claimed spectrum. Apply
Lemma~\ref{lem:symmetry} to obtain its stationary marginal assertion.
For the uniform weight, the Perron vector and stationary law are uniform,
and~\eqref{eq:product} supplies the comparison chain.
\end{proof}

\section{An exact parity control with different dynamics}
The positive parity family
\begin{equation}\label{eq:parity}
 \mu_t(x)=2^{-L}(1+t\chi_{[L]}(x)),\qquad |t|<1,
 \qquad \chi_S(x)=(-1)^{\sum_{i\in S}x_i},
\end{equation}
has all proper marginals equal to product marginals. Its single-site
heat-bath gap tends to zero as $|t|\to1$ at fixed size $L\ge2$ and fixed
finite positive site rates~\cite{ErikssonParity}. Nevertheless, its transfer ratio at fixed
$\beta$ is bounded uniformly away from one, as follows.

\begin{proposition}\label{prop:parity}
Let $L\ge2$, $\rho=\tanh\beta$, and $|t|<1$. Up to a common positive
scalar, the spectrum of $A_{\mu_t}$ is the union over unordered complement
pairs $\{S,S^c\}$ of the two roots
\begin{equation}\label{eq:parityroots}
 \frac{z_s\pm\sqrt{z_s^2-4\Delta}}2,
 \quad z_s=\rho^s+\rho^{L-s},\quad s=|S|,
 \quad \Delta=(1-t^2)\rho^L.
\end{equation}
The Perron root is the upper root at $s=0$, and the second eigenvalue is
the upper root at $s=1$. In particular,
\begin{equation}\label{eq:paritybound}
 \frac{\lambda_2(A_{\mu_t})}{\lambda_1(A_{\mu_t})}
 \le\frac{\rho+\rho^{L-1}}{1+\rho^L}
 \le\tanh(2\beta)<1.
\end{equation}
The last bound is sharp as a supremum over $L\ge2$ and $|t|<1$.
\end{proposition}
\begin{proof}
The row-normalized product transfer $P_\beta=M_\beta/(1+q)^L$ acts on
$\chi_S$ by $\rho^{|S|}$. Write $D_t=\diag\sqrt{1+t\chi_{[L]}}$.
The weighted congruence, apart from a common scalar, is $D_tP_\beta D_t$,
which is similar to $P_\beta D_t^2$ since $D_t$ is invertible.
Multiplication by $\chi_{[L]}$ exchanges $\chi_S$ and $\chi_{S^c}$.
The coefficient block on that ordered pair is
\[
 \begin{pmatrix}
 \rho^s&t\rho^s\\ t\rho^{L-s}&\rho^{L-s}
 \end{pmatrix}.
\]
Its trace is $z_s$ and its determinant is $\Delta$, yielding
\eqref{eq:parityroots}. The $2^{L-1}$ complement pairs exhaust the
character basis, with repetitions counted as usual.

The largest $z_s$ is $z_0=1+\rho^L$, and the largest among proper pairs
is $z_1=\rho+\rho^{L-1}$. The upper root increases with $z_s$.
The lower root in the empty/full pair is at most $\rho^L$, whereas
the upper root for $s=1$ is at least $\rho$. Every other lower root
is no larger than its own upper root. This identifies the two largest
eigenvalues.

For fixed $\Delta>0$, the upper root divided by $z$ is
\[
 \frac{1+\sqrt{1-4\Delta/z^2}}2,
\]
which increases with $z\ge2\sqrt\Delta$. Comparing $z_1\le z_0$ proves
the first inequality in~\eqref{eq:paritybound}. For the second, write
$b=\rho^{L-2}\le1$; the ratio is $\rho(1+b)/(1+\rho^2b)$, increasing
in $b$. Its maximum is $2\rho/(1+\rho^2)=\tanh(2\beta)$.
For $L=2$ and $t\to1$, $\Delta\to0$, so equality is approached.
\end{proof}

This control rules out a tempting shortcut: the known parity heat-bath
obstruction cannot be transferred by changing the name of the operator.
The code construction uses a different geometric ingredient, a set of
exponentially many mutually separated configurations while retaining
surjective coordinate projections.

\section{What marginal certificates cannot establish}
\begin{corollary}\label{cor:certificate}
Fix the constants of Theorem~\ref{thm:main}. Suppose $F$ is a valid lower
bound on the transfer gap for every strictly positive weight, and factors
through the tuple
\[
 (L,\beta,\mathcal M_{r_L}(\mu),\mathcal M_{r_L}(\pi)).
\]
At the product input and for all sufficiently large $L$, its value is at
most $7e^{-\gamma L}$. The same conclusion holds if only one of the two
marginal tuples is supplied.
\end{corollary}
\begin{proof}
At the same $L$ and $\beta$, the constructed weight and its stationary
law supply exactly the same input as the uniform weight and uniform
stationary law. A function of that input takes the same value on both.
Validity on the constructed example and~\eqref{eq:slow} for $j=2$
force the asserted bound.
\end{proof}

The corollary does not assert that every such functional is zero or that
the product chain is slow. Its actual gap is $1-\tanh\beta$. Nor does
it concern a certificate using full conditionals, uniform density bounds,
the operator entries, or code membership: those are additional information.
In particular, the weights here can have exponentially small probabilities
and global many-body interactions. The theorem says nothing against a
result restricted to uniformly controlled, fixed-range Gibbs weights.
The Perron Doob chain updates a full configuration and is not a local
single-site dynamics. We make no claim about continuum Yang--Mills,
a physical mass gap, entropy tensorization, or a specified spatial mixing
rate. The many-mode conclusion is spectral; no separate theorem about
metastable sets is implicit.

\section{Upper density control without a lower bound}\label{sec:density}
The obstruction above allows the maximum density relative to product to
increase with dimension. A uniform upper bound changes the conclusion;
a positive lower bound is not required. Write $\rho=\tanh\beta$ and let
$P_\rho$ be the Boolean noise operator in $L^2(u_L)$, whose Walsh
multipliers are $\rho^{|S|}$. For a nonnegative density $r$ with
$\mathbb E_{u_L}r=1$, put
\[
 T_r=\diag(\sqrt r)P_\rho\diag(\sqrt r).
\]
Since $P_\rho=M_\beta/(1+e^{-2\beta})^L$, for $\mu=2^{-L}r$ we have
$A_\mu=2^{-L}(1+e^{-2\beta})^L T_r$. Thus the eigenvalue ratios agree.
Here zeros of $r$ are allowed; we keep the zero eigenvalues outside its
support in the full $2^L$-dimensional list.

\begin{theorem}\label{thm:density}
Let $L\ge1$, $0<\rho<1$, $B\ge1$, and $0\le r\le B$ with
$\mathbb E_{u_L}r=1$. Set
\[
 e_\rho=\frac{1+\rho^2}{(1-\rho)^2},\qquad
 c(B,\rho)=B^{-1}(2B)^{-2e_\rho}.
\]
The two largest eigenvalues of $T_r$ satisfy
$1-\lambda_2/\lambda_1\ge c(B,\rho)>0$, independently of $L$.
In particular, no positive lower bound on $r$ is needed.
\end{theorem}
\begin{proof}
All norms and scalar products use $u_L$. The noise operator is positive
semidefinite, fixes constants and is a contraction. Hence
$\lambda_1\le B$. The unit vector $\sqrt r$ has Rayleigh quotient
$\langle r,P_\rho r\rangle\ge(\mathbb E r)^2=1$, by Walsh expansion,
so $1\le\lambda_1\le B$. On $S=\{r>0\}$ the matrix is strictly positive;
its Perron vector is positive on $S$ and its leading eigenvalue is simple.

We use two classical Boolean hypercontractive inequalities. Put
$q_*=1+\rho^{-2}$ and $p_*=1-\rho$. Then
\begin{equation}\label{eq:hyperinputs}
 \|P_\rho f\|_{q_*}\le\|f\|_2,\qquad
 \langle f,P_\rho g\rangle\ge\|f\|_{p_*}\|g\|_{p_*}\quad(f,g\ge0).
\end{equation}
For clarity about conventions, equations (1.1)--(1.2) and tensorization
in~\cite{MOS} use time $t=-\log\rho$ for the uniform Boolean semigroup.
For the second inequality set $p_*'=p_*/(p_*-1)=-(1-\rho)/\rho$.
Then
$\tfrac12\log((1-p_*')/(1-p_*))=-\log\rho$.
The reverse inequality gives $\|P_\rho g\|_{p_*'}\ge\|g\|_{p_*}$;
reverse H\"older gives the displayed bilinear bound. This is first applied
to positive functions, then to $f+\varepsilon,g+\varepsilon$ and the limit
$\varepsilon\downarrow0$. Finite sums and $p_*>0$ justify that limit;
we do not take a negative moment of a function with zeros.

If $\lambda_2<\lambda_1/2$, the gap exceeds $1/2$. Otherwise
$\lambda_2\ge1/2$. Choose a real unit second eigenvector $w$ orthogonal
to the Perron vector. It vanishes off $S$ and has both signs. Write
$w=s h_+-t h_-$, where $s,t>0$, $s^2+t^2=1$, and $h_+,h_-$ are
nonnegative orthonormal functions with disjoint supports. Positivity
and the eigenrelation give pointwise
$\lambda_2h_\pm\le T_rh_\pm$: on the corresponding sign support this
follows from the eigenrelation, and elsewhere from positivity.
Integration against the unit $h_\pm$ gives
$\lambda_2\le\langle\sqrt r h_\pm,P_\rho\sqrt r h_\pm\rangle$.
For $f_\pm=\sqrt r\,h_\pm$ this yields
\[
 \lambda_2f_\pm\le rP_\rho f_\pm,\qquad
 \|f_\pm\|_{q_*}\le\frac B{\lambda_2}\|f_\pm\|_2,\qquad
 \|f_\pm\|_2^2\ge\langle f_\pm,P_\rho f_\pm\rangle\ge\lambda_2.
\]
The middle inequality uses the first part of~\eqref{eq:hyperinputs}.
No division by $r$ occurs.

Let $\theta=p_*(q_*-2)/(2(q_*-p_*))$. Ordinary H\"older applied to
powers of $f$ gives
$\|f\|_2\le\|f\|_{p_*}^{\theta}\|f\|_{q_*}^{1-\theta}$,
since $1/2=\theta/p_*+(1-\theta)/q_*$. This does not require a normed
space at exponent $p_*<1$: the H\"older weights
$2\theta/p_*$ and $2(1-\theta)/q_*$ are positive and sum to one.
Thus each sign part satisfies
\[
 \|f_\pm\|_{p_*}\ge\sqrt{\lambda_2}
       (B/\lambda_2)^{-e_\rho},\qquad
 \frac{1-\theta}{\theta}
 =\frac{q_*(2-p_*)}{p_*(q_*-2)}=e_\rho.
\]
By the second part of~\eqref{eq:hyperinputs}, the cross entry obeys
\[
 d=\langle h_+,T_rh_-\rangle
   =\langle f_+,P_\rho f_-\rangle
   \ge\lambda_2(B/\lambda_2)^{-2e_\rho}.
\]
The compression of $T_r$ to the basis $(h_+,h_-)$ is
$\left(\begin{smallmatrix}a&d\\d&b\end{smallmatrix}\right)$.
Its eigenvector $(s,-t)$ has eigenvalue $\lambda_2$, giving
$a-\lambda_2=dt/s$ and $b-\lambda_2=ds/t$.
Its other eigenvalue is $\lambda_2+d/(st)$, because $s^2+t^2=1$.
The variational principle and $st\le1/2$ now imply
$\lambda_1-\lambda_2\ge d/(st)\ge2d$. Consequently
\[
 1-\frac{\lambda_2}{\lambda_1}
 \ge\frac{2\lambda_2}{\lambda_1}(B/\lambda_2)^{-2e_\rho}
 \ge B^{-1}(2B)^{-2e_\rho}.
\]
This does not deteriorate when one sign part has small norm. In the
previously omitted case, $c(B,\rho)\le1/4<1/2$, since $e_\rho\ge1$.
If the support has one point, $\lambda_2=0$ and the gap is one, so that
case is included as well.
\end{proof}

\begin{corollary}\label{cor:densitygrowth}
At fixed $\rho$, suppose a sequence has
$\gap(T_{r_L})\le K e^{-\gamma L}$, with $K,\gamma>0$. Then
\[
 \max r_L\ge
 \left(\frac{2^{-2e_\rho}}K\right)^{1/(2e_\rho+1)}
       e^{\gamma L/(2e_\rho+1)}.
\]
\end{corollary}
\begin{proof}
Apply Theorem~\ref{thm:density} with $B=\max r_L\ge1$ and rearrange.
\end{proof}
For Theorem~\ref{thm:main}, take $K=7$ and its positive $\gamma$.
The code mixture has maximum density $2^L(1-\eta)/|C|+\eta$, consistent
with this necessary growth. Upper density growth alone is not asserted
to be sufficient for collapse. Tiny minimum density alone cannot cause
collapse under a uniform maximum-density bound. These are statements
about this transfer family, not single-site heat-bath dynamics.
The constants are not claimed optimal.

\section{Finite density moments: a sufficient regime and an obstruction}
\label{sec:finite-moment}
A pointwise density cap is stronger than necessary. We next replace it
by a dimension-independent finite moment. A second construction shows why
one cannot simply assume that any finite moment suffices. All norms in
this section use the uniform probability law $u_L$; the operator is the
same $T_r$ as in Section~\ref{sec:density}.

\begin{theorem}\label{thm:moment}
Fix $0<\rho<1$, $s>(1+\rho)/(1-\rho)$ and $B\ge1$. For every $L\ge1$
and density $r\ge0$ with $\mathbb E r=1$ and $\|r\|_s\le B$,
\[
 \gap(T_r)\ge\min\{1/2,(2B)^{-2e}\}>0,
 \qquad e=\frac{2\rho s}{(1-\rho)(s(1-\rho)-1-\rho)}.
\]
Neither a pointwise upper bound nor a positive lower bound is assumed.
\end{theorem}
\begin{proof}
Set $A=1+\rho$, $p=1-\rho$, $t=1+1/\rho$ and
\[
 Q=\frac{st}{s+t}=\frac{s(1+\rho)}{s\rho+1+\rho}.
\]
The condition on $s$ is exactly $Q>A$, so $Q>A>1>p>0$.
Ordinary Boolean hypercontractivity gives
\[
 \|P_{\sqrt\rho}f\|_2\le\|f\|_A,
 \qquad \|P_\rho f\|_t\le\|f\|_A.
\]
Indeed the noise thresholds are respectively
$\sqrt{(A-1)/(2-1)}=\sqrt\rho$ and
$\sqrt{(A-1)/(t-1)}=\rho$. The reverse bilinear inequality in
\eqref{eq:hyperinputs} applies at exponent $p$.

As before, testing the unit vector $\sqrt r$ gives $\lambda_1\ge1$.
For any real unit vector $h$, put $f=\sqrt r\,h$ and
$A_0=2s/(s+1)>A$. H\"older and probability-space monotonicity give
\[
 \langle h,T_rh\rangle=\|P_{\sqrt\rho}f\|_2^2
 \le\|f\|_A^2\le\|f\|_{A_0}^2
 \le\|\sqrt r\|_{2s}^2\|h\|_2^2\le B,
\]
where $1/A_0=1/(2s)+1/2$. Thus $1\le\lambda_1\le B$ without
using a maximum-density estimate.

If $\lambda_2<\lambda_1/2$, the gap is greater than $1/2$.
Otherwise $\lambda_2\ge1/2$. On the nonempty support of $r$ the matrix
has strictly positive entries. A real unit second eigenvector vanishes
off this support and has both signs, by orthogonality to the positive
Perron vector. Write it as $w=c h_+-d h_-$ with $c,d>0$,
$c^2+d^2=1$ and nonnegative orthonormal disjoint sign parts.
Positivity and the eigenrelation imply
$\lambda_2h_\pm\le T_rh_\pm$ pointwise, including their zero sets.
For either $f=\sqrt r\,h_\pm$, multiplication by $\sqrt r$ gives
$\lambda_2f\le rP_\rho f$. H\"older with $1/Q=1/s+1/t$ yields
\[
 \|f\|_Q\le (B/\lambda_2)\|f\|_A.
\]
Integration against the unit sign part also gives
$\lambda_2\le\langle f,P_\rho f\rangle\le\|f\|_A^2$.
Let $0<\theta<1$ solve $1/A=\theta/p+(1-\theta)/Q$.
H\"older applied to powers, with positive weights
$A\theta/p$ and $A(1-\theta)/Q$ summing to one, gives
$\|f\|_A\le\|f\|_p^\theta\|f\|_Q^{1-\theta}$.
No triangle inequality at exponent $p<1$ is used. Combining these bounds,
\[
 \|f\|_p\ge\sqrt{\lambda_2}(B/\lambda_2)^{-e},\qquad
 e=\frac{1-\theta}{\theta}
   =\frac{Q(A-p)}{p(Q-A)}
   =\frac{2\rho s}{(1-\rho)(s(1-\rho)-1-\rho)}.
\]
Reverse hypercontractivity now bounds the cross entry by
$D=\langle h_+,T_rh_-\rangle\ge\lambda_2(B/\lambda_2)^{-2e}$.
The two-dimensional compression, exactly as in the proof of
Theorem~\ref{thm:density}, has eigenvalues $\lambda_2$ and
$\lambda_2+D/(cd)$. Since $cd\le1/2$, the variational principle gives
\[
 \gap(T_r)\ge\frac{2D}{\lambda_1}
 \ge\frac{2\lambda_2}{\lambda_1}(B/\lambda_2)^{-2e}
 \ge(2B)^{-2e}.
\]
Here $2\lambda_2/\lambda_1\ge1$ and $\lambda_2\ge1/2$.
Combining the two cases proves the claim. One-point support has
$\lambda_2=0$ and is already covered. No step divides by $r$.
\end{proof}

\begin{theorem}\label{thm:spikes}
Fix $0<\rho<1$, $B>1$ and
\[
 1<s<s_{\rm point}:=\frac{\log2}{\log(2/(1+\rho))}.
\]
There are strictly positive normalized densities $r_L$ for every $L\ge1$
with $\|r_L\|_s\le B$ whose actual transfer gaps tend to zero
exponentially. Explicitly put $\varepsilon=(B-1)/2^{1/s}$ and
\[
 z_L=\varepsilon2^{L/s},\qquad
 r_L(x)=\frac{1+z_L(\mathbf1_{\{0^L\}}(x)+\mathbf1_{\{1^L\}}(x))}{Z_L},
 \qquad Z_L=1+2\varepsilon2^{-L(1-1/s)}.
\]
Write $a=(1+\rho)/2$, $b=(1-\rho)/2$, $v=(1+\rho^2)/2$ and
\[
 \kappa=2^{1/s}a,\qquad
 \eta=2^{-1/(2s)}\frac{\sqrt v}{a},\qquad \zeta=b/a.
\]
Then $\kappa>1$, $0<\eta,\zeta<1$, and
\begin{equation}\label{eq:spikegap}
 \gap(T_{r_L})\le2\zeta^L+\frac{1+2\varepsilon}{\varepsilon}\kappa^{-L}
                      +\sqrt{2/\varepsilon}\,\eta^L.
\end{equation}
\end{theorem}
\begin{proof}
The antipodes are distinct for every $L\ge1$. The formula for $Z_L$
is their exact normalization, so $\mathbb E r_L=1$ and $r_L>0$.
The $L^s$ triangle inequality gives
\[
 \|1+z_L(\mathbf1_{\{0^L\}}+\mathbf1_{\{1^L\}})\|_s
 \le1+z_L(2/2^L)^{1/s}=B.
\]
Division by $Z_L\ge1$ proves the moment bound for all $L$.

Use the Euclidean matrix of $T_{r_L}$; the common normalization of the
$L^2(u_L)$ coordinate basis leaves this matrix unchanged. Split the two
antipodes $C$ from their complement. The $C$ principal block is
\[
 D_L\begin{pmatrix}1&\zeta^L\\\zeta^L&1\end{pmatrix},
 \qquad D_L=(1+z_L)a^L/Z_L.
\]
Principal interlacing gives $\lambda_2\ge D_L(1-\zeta^L)$, and a
coordinate Rayleigh vector gives $\lambda_1\ge D_L$. The complementary
block is $Z_L^{-1}$ times a principal submatrix of the positive
contraction $P_\rho$, so its norm is at most one. For the rectangular
cross block $K$, each full squared row sum of $P_\rho$ is
$(a^2+b^2)^L=v^L$. Therefore
\[
 \|K\|^2\le\|K\|_{\rm F}^2\le2(1+z_L)v^L/Z_L^2.
\]
The symmetric off-diagonal block has norm $\|K\|$. Consequently
$\lambda_1\le D_L(1+\zeta^L)+1+\|K\|$, which implies
\[
 \gap(T_{r_L})\le2\zeta^L+D_L^{-1}+\|K\|/D_L.
\]
Since $Z_L\le1+2\varepsilon$ and $1+z_L\ge\varepsilon2^{L/s}$,
\[
 D_L^{-1}\le\frac{1+2\varepsilon}{\varepsilon}\kappa^{-L},
 \qquad \frac{\|K\|}{D_L}\le\sqrt{2/\varepsilon}\,\eta^L.
\]
This proves~\eqref{eq:spikegap}. The strict condition on $s$ says
$\kappa>1$. Also $\eta^2=v/(a\kappa)<v/a<1$, since $\rho^2<\rho$.
Thus all three terms decay exponentially. At $L=1$ the complementary
block is empty; the same bounds remain valid. The estimate need not be
informative at small $L$.
\end{proof}

These regimes leave a genuine interval open. With
$t_0=(1-\rho)/(1+\rho)\in(0,1)$, strict concavity gives
$\log(1+t_0)>t_0\log2$, hence
\[
 1<s_{\rm point}<s_c:=\frac{1+\rho}{1-\rho}.
\]
Neither the intermediate interval nor either endpoint is settled here;
no sharp threshold is claimed. At $B=1$, Jensen forces $r=1$ and the
actual gap is $1-\rho$, displaying the weakness of the sufficient bound.
The spike family does not claim product marginals for its weighting or
stationary law. It isolates a finite-moment obstruction for the same
transfer operator, whereas Theorem~\ref{thm:main} isolates the distinct
marginal-information obstruction. Both are compatible with
Theorems~\ref{thm:density} and~\ref{thm:moment}.

\section{Comparison with hyperbounded and perturbation results}
The inequalities and matrix arguments used above are classical. Related
literature already supplies both qualitative gap criteria and, under
different hypotheses, dimension-independent quantitative perturbation
bounds. Miclo's Theorem~1 in~\cite{Miclo} gives a gap for ergodic
self-adjoint hyperbounded Markov operators. Its Proposition~11 explains
why a general quantitative gap does not follow from hyperboundedness
alone. Wang's Theorem~1.1 in~\cite{Wang} gives a tail-norm criterion
for the gap of an invariant Markov operator's symmetrization.
These fixed-space conclusions must not be confused with a bound uniform
over a family of increasing dimensions.

Here $T_r$ is generally not Markov in $u_L$. Its Perron Doob transform
is Markov in its own stationary law, which generally differs from $r u_L$.
The ground-state multiplication is an $L^2$ isometry between these spaces,
but does not automatically identify their higher norms uniformly in $L$.
Thus an implication from a qualitative hyperboundedness theorem would
still need a quantitative change-of-measure argument.

Gross's Theorem~2.2 in~\cite{Gross} does give dimension-independent gaps
for additive Schr\"odinger perturbations on Riemannian probability
manifolds with a base logarithmic Sobolev inequality and suitable moments
of both $e^V$ and $e^{-V}$. We have not identified that additive operator
with the logarithm of this discrete transfer congruence. Our moment
hypothesis supplies no inverse-density moment and permits zeros.
Gl\"uck~\cite{Gluck} gives hyperboundedness criteria and uniform essential
spectral-radius estimates; those estimates alone do not control the
second eigenvalue of a finite matrix. These differences do not prove
originality or exclude a stronger implication. The retained comparison
also records incomplete access to Aida's original gap theorem.
No external priority, improvement over this literature, or exhaustive
novelty clearance is asserted.

\section{Relation to predecessors and reproducibility}
The congruence analysis in~\cite{ErikssonCongruence} permits positive
weights concentrating on an antipodal pair, with a transfer ratio tending
to one as volume grows. Such a pair fails to preserve pair marginals.
The parity analysis in~\cite{ErikssonParity} preserves every proper
marginal but concerns a heat-bath generator; Proposition~\ref{prop:parity}
shows precisely why it does not supply the transfer conclusion.
The present synthesis obtains linear-order marginal matching for both
the weight and the Perron stationary law, together with exponentially
many slow transfer modes.

The code ingredients and the matrix methods are classical. In particular,
Lemma~\ref{lem:code} is an elementary existence argument, not a new code
construction and not an efficient encoding claim. The citation to~\cite{AGM} supplies general limited-independence background, not a direct historical attribution for the dual-distance projection equivalence proved here. Work on simultaneously
good primal and dual codes, including~\cite{BrehmResch}, is stronger in
algorithmic aspects. Our targeted primary-source check did not identify
the exact joint conclusion above; this is not proof of historical priority.
External originality and the importance of this specialization remain
unestablished.

All proofs are finite-dimensional and analytical. No numerical experiment
or formal verification is needed or claimed. A source package retains the
preregistered proof checks, complete analytical audit, source hashes and
this manuscript. Document reproduction is a separate check of exact bytes
and layout, not an independent proof of the theorems.

\section*{Model assessments and review context}
OpenAI/Codex assisted the mathematical synthesis and manuscript preparation.
The fabrication session is involved in the work. A distinct involved
Codex/OpenAI session audits the frozen mathematical source; exact served
backend versions and reasoning settings are not exposed, so no particular
model version or provider independence is asserted. Analytical source
review is not human peer review and supplies no inherited PDF rating.
Any assessment attached to this manuscript must identify the exact final
PDF hash, its actual review context, complete report, scores and objections.
The manuscript does not incorporate a grade from another version or paper.
The author owns and administers AIRR. A working-paper deposit is not an
editorial acceptance or evidence of human peer review.

\smallskip\noindent
Reproducible source commit:\par\noindent
{\small\ttfamily\SourceCommit}

\begin{thebibliography}{99}
\bibitem{AGM}
N. Alon, O. Goldreich and Y. Mansour,
Almost $k$-wise independence versus $k$-wise independence,
\emph{Information Processing Letters} \textbf{88} (2003), 107--110.
\href{https://www.wisdom.weizmann.ac.il/~oded/PSX/agm.pdf}{Author manuscript}.
\bibitem{BrehmResch}
M. Brehm and N. Resch,
Linear Time Encodable Binary Code Achieving GV Bound with Linear Time
Encodable Dual Achieving GV Bound,
\emph{ITCS 2026}, LIPIcs \textbf{362}, article 28.
\href{https://doi.org/10.4230/LIPIcs.ITCS.2026.28}{Publisher version}.
\bibitem{ErikssonCongruence}
L. Eriksson,
Congruence Rigidity and the Fusion Bound: What a Positive Weight Can and
Cannot do to the Spectrum of a Transfer Kernel,
AIRR, ARR-2026-7S1XA3PZEB9QM8FQ, version 1.
\href{https://airr.science/papers/ARR-2026-7S1XA3PZEB9QM8FQ/versions/v1/}{Public record}.
\bibitem{ErikssonParity}
L. Eriksson,
Strictly Positive Parity Laws and the Limits of Marginal Certificates:
Exact Heat-Bath Spectrum and Optimal Variance Tensorization,
AIRR, ARR-2026-5HRC3PS9HZ9WDAA5, version 2 (2026).
\href{https://airr.science/papers/ARR-2026-5HRC3PS9HZ9WDAA5/versions/v2/}{Public working paper}.
\bibitem{MOS}
E. Mossel, K. Oleszkiewicz and A. Sen,
On Reverse Hypercontractivity, arXiv:1108.1210v5 (2012).
\href{https://arxiv.org/abs/1108.1210}{Primary manuscript}.
\bibitem{Gross}
L. Gross, Invariance of intrinsic hypercontractivity under perturbation of
Schr\"odinger operators, arXiv:2412.20282v2 (2025).
\href{https://arxiv.org/abs/2412.20282}{Primary manuscript}.
\bibitem{Gluck}
J. Gl\"uck, Spectral gaps for hyperbounded operators,
arXiv:1802.09422v6 (2019).
\href{https://arxiv.org/abs/1802.09422}{Primary manuscript}.
\bibitem{Miclo}
L. Miclo, On hyperboundedness and spectrum of Markov operators.
\href{https://perso.math.univ-toulouse.fr/miclo/files/2012/04/hyperboundedness.pdf}{Author manuscript}.
\bibitem{Wang}
F.-Y. Wang, Criteria of Spectral Gap for Markov Operators,
arXiv:1305.4460v6 (2013).
\href{https://arxiv.org/abs/1305.4460v6}{Primary manuscript}.
\end{thebibliography}
\end{document}